\section{Фазовые переходы}
\subsection{Удельная теплота сгорания}
%%%%%%%%%%%%%%%%%%%%%%%%%%%%%%%%%%%%%%%%%%%%%%%%%%%%
%%%%%%%%%%%%%%%%%%%%%%%%%%%%%%%%%%%%%%%%%%%%%%%%%%%%
 \z{Медный сосуд массой $m_{\text{с}} = 500\,\text{г}$ содержит $V = 2\,\text{л}$ воды при температуре $t_1 = 10\,\degree\text{C}$. До какой температуры можно нагреть воду, сжигая спирт массой $m_{\text{сп}} = 50\,\text{г}$? КПД горелки $\eta = 50\%$.}
%
{$T_{\text{к}} = T_0 + \frac{m_{\text{спирт}} \cdot q \cdot \eta}{(m_{\text{в}} c_{\text{в}} + m_{\text{м}} c_{\text{м}})} = 88.5\,\degree\text{C}$}

%%%%%%%%%%%%%%%%%%%%%%%%%%%%%%%%%%%%%%%%%%%%%%%%%%%%
 \z{При помощи нагревателя с КПД $\eta = 40\%$ необходимо довести до температуры кипения $V = 4\,\text{л}$ воды в алюминиевой кастрюле массой $m_{\text{к}} = 2\,\text{кг}$. Определите расход керосина на нагревание воды и кастрюли, если их начальная температура $t_1 = 20\,\degree\text{C}$.}
%
{$m_{\text{к}} = \frac{(m_{\text{в}} c_{\text{в}} + m_{\text{а}} c_{\text{а}}) \Delta T}{q \eta} = 0.08\,\text{кг}$}

%%%%%%%%%%%%%%%%%%%%%%%%%%%%%%%%%%%%%%%%%%%%%%%%%%%%
 \z{В алюминиевом чайнике массой $m_{\text{ч}} = 900\,\text{г}$ нагрели $V = 10\,\text{л}$ воды до кипения, израсходовав природный газ массой $m_{\text{г}} = 200\,\text{г}$. Определите начальную температуру воды. КПД горелки $\eta = 40\%$.}
%
{$T_0 = T_{100} - \frac{m_{\text{г}} \cdot q \cdot \eta}{m_{\text{в}} c_{\text{в}} + m_{\text{а}} c_{\text{а}}} = 18\,\degree\text{C}$}

%%%%%%%%%%%%%%%%%%%%%%%%%%%%%%%%%%%%%%%%%%%%%%%%%%%%
 \z{В медной кастрюле нагрели $V = 5\,\text{л}$ воды от температуры $t_1 = 10\,\degree\text{C}$ до кипения, израсходовав керосин массой $m_{\text{кер}} = 100\,\text{г}$. Определите массу кастрюли, если КПД нагревателя $\eta = 40\%$.}
%
{$m_{\text{к}} = \frac{m_{\text{кер}} q \eta}{\Delta T c_{\text{м}}} - m_{\text{в}} \frac{c_{\text{в}}}{c_{\text{м}}} = -5\,\text{кг}$}

%%%%%%%%%%%%%%%%%%%%%%%%%%%%%%%%%%%%%%%%%%%%%%%%%%%%
 \z{На газовой плите за $\tau = 15\,\text{мин}$ довели до температуры кипения $V = 3\,\text{л}$ воды. Какова масса газа, сгоравшего каждую секунду, если начальная температура воды $t_1 = 20\,\degree\text{C}$?}
%
{$\frac{\Delta m_{\text{г}}}{\Delta t} = \frac{c_{\text{в}} m_{\text{в}} \Delta T}{q \tau} = 0.026\,\text{г/с}$}

%%%%%%%%%%%%%%%%%%%%%%%%%%%%%%%%%%%%%%%%%%%%%%%%%%%%
 \z{Мощность автомобильного двигателя $P = 44\,\text{кВт}$. Каков КПД этого двигателя, если при скорости $v = 100\,\text{км/ч}$ он на $s = 100\,\text{км}$ пути расходует бензин объемом $V = 14\,\text{л}$?}
%
{$\eta = \frac{3600 \cdot 44 \cdot 10^3}{14 \cdot 44 \cdot 10^6 \cdot 0.74} = 35\%$}

%%%%%%%%%%%%%%%%%%%%%%%%%%%%%%%%%%%%%%%%%%%%%%%%%%%%
 \z{Междугородный автобус проехал $s = 160\,\text{км}$ за $t = 2\,\text{ч}$, развивая при этом мощность $P = 70\,\text{кВт}$. Сколько литров бензина израсходовал автобус, если КПД его двигателя $\eta = 25\%$?}
%
{$m_{\text{б}} = \frac{P \tau}{q \eta} = 45\,\text{кг} = 62\,\text{л}$}

